\clearpage\pagestyle{otMain}\fontsize{11}{15}\selectfont\corrigesChapitre{18}{Raisonner et préparer le brevet}
\begin{blocLivre}{cor-18-1}\begin{corrige}{1}{Un contre-exemple}\demarche{Chercher un cas qui contredit exactement l’affirmation.}
Avec $x=1/2$, $x^2=1/4<1/2$. Ce seul contre-exemple suffit. Avec 0 ou 1, l’inégalité stricte serait également fausse par égalité.
\end{corrige}\end{blocLivre}
\begin{blocLivre}{cor-18-2}\begin{corrige}{2}{Deux entiers impairs}\demarche{Traduire l’hypothèse par une écriture générale.}
Écrivons les deux impairs $2a+1$ et $2b+1$, avec a, b entiers. Leur somme vaut $2(a+b+1)$, donc c’est un multiple de 2. La preuve vaut pour tous les entiers concernés.
\end{corrige}\end{blocLivre}
\begin{blocLivre}{cor-18-3}\begin{corrige}{3}{Contrôler une unité}\demarche{Convertir les longueurs avant le produit.}
40 cm valent 0,4 m. L’aire vaut $3\times0{,}4=\res{1{,}2\text{ m}^2}$. Le calcul 120 a multiplié des nombres exprimés dans des unités différentes.
\end{corrige}\end{blocLivre}
\begin{blocLivre}{cor-18-4}\begin{corrige}{4}{Un arrondi utile}\demarche{Interpréter la quantité entière dans le contexte.}
$7{,}2/3=2{,}4$. Deux bidons n’apportent que 6 L ; trois apportent 9 L. Il faut $\res3$ bidons. L’arrondi à l’entier le plus proche ne couvre pas forcément le besoin.
\end{corrige}\end{blocLivre}
\begin{blocLivre}{cor-18-5}\begin{corrige}{5}{Une hypothèse manquante}\demarche{Identifier l’hypothèse nécessaire.}
Non : il faudrait savoir que les côtés de 3 et 4 cm sont perpendiculaires pour appliquer Pythagore ainsi. Un triangle de côtés 3, 4, 6 existe aussi car $3+4>6$, et son troisième côté n’est pas 5.
\end{corrige}\end{blocLivre}
\begin{blocLivre}{cor-18-6}\begin{corrige}{6}{Un programme constant}\demarche{Passer des essais à l’expression générale.}
Le résultat est $2(x+5)-2x=2x+10-2x=\res{10}$. Des essais suggèrent ce résultat ; la simplification le prouve pour tout x.
\end{corrige}\end{blocLivre}
\begin{blocLivre}{cor-18-7}\begin{corrige}{7}{Comparer exactement}\demarche{Conserver une comparaison exacte.}
Les dénominateurs sont positifs. $17\times7=119$ et $5\times24=120$, donc $\res{17/24<5/7}$. Un arrondi au dixième aurait masqué la différence.
\end{corrige}\end{blocLivre}
\begin{blocLivre}{cor-18-8}\begin{corrige}{8}{Une conjecture de divisibilité}\demarche{Distinguer les deux parités possibles.}
Deux entiers consécutifs comprennent toujours un pair : si n est pair, le produit contient ce facteur ; sinon n+1 est pair. Dans les deux cas, $n(n+1)$ est multiple de 2.
\end{corrige}\end{blocLivre}
\begin{blocLivre}{cor-18-9}\begin{corrige}{9}{Vérifier une équation}\demarche{Une substitution rapide détecte une erreur.}
Pour 4, les membres valent 6 et 8 : ce n’est pas une solution. $3x-6=x+4$ donne $2x=10$, donc $\res{x=5}$. Le contrôle donne 9 de chaque côté.
\end{corrige}\end{blocLivre}
\begin{blocLivre}{cor-18-10}\begin{corrige}{10}{Une moyenne mal interprétée}\demarche{Construire un contre-exemple satisfaisant la donnée.}
Non. Considérons quatre notes $10;10;10;18$ : leur moyenne vaut 12, mais une seule est au moins 12. La propriété de partage en deux moitiés concerne la médiane, pas la moyenne.
\end{corrige}\end{blocLivre}
\begin{blocLivre}{cor-18-11}\begin{corrige}{11}{Une commande au meilleur prix}\demarche{Inclure conditionnement et frais fixes.}
A coûte $24\times2{,}80+6=\res{73{,}20\text{ €}}$. Il faut trois lots chez B : $\res{63\text{ €}}$. B économise $\res{10{,}20\text{ €}}$. On compare les coûts totaux pour la même quantité.
\end{corrige}\end{blocLivre}
\begin{blocLivre}{cor-18-12}\begin{corrige}{12}{Un modèle de bassin}\demarche{Relier volume, aire de base et hauteur.}
6000 L valent 6 m$^3$. La base vaut 10 m$^2$, donc la hauteur est $6/10=\res{0{,}6\text{ m}}$. Le volume maximal vaut 8 m$^3$ ; il reste $\res{2000\text{ L}}$.
\label{fig-cor-18-12}\figureGeo[corrige]{synthese-bassin}
\end{corrige}\end{blocLivre}
\begin{blocLivre}{cor-18-13}\begin{corrige}{13}{Un hasard et un budget}\demarche{Séparer modèle aléatoire et bilan comptable.}
La probabilité de gain vaut $20/200=\res{1/10}$. Les recettes sont 600 € et les lots coûtent 300 €. Le solde vaut $\res{300\text{ €}}$. Le résultat financier global ne décrit pas la chance d’un billet particulier.
\end{corrige}\end{blocLivre}
\begin{blocLivre}{cor-18-14}\begin{corrige}{14}{Un maximum expliqué}\demarche{Transformer l’expression pour lire une borne atteinte.}
L’autre côté vaut $20-x$, avec $0<x<20$. L’aire vaut $20x-x^2=100-(x-10)^2$. Un carré étant positif ou nul, l’aire est au plus $\res{100\text{ cm}^2}$. Cette valeur est atteinte pour $x=10$, donc pour un carré.
\label{fig-cor-18-14}\figureGeo[corrige]{rectangle-fonction}
\end{corrige}\end{blocLivre}
\begin{blocLivre}{cor-18-15}\begin{corrige}{15}{Un plan de révision}\demarche{Relier le bilan à une action précise.}
Les automatismes rapportent $8\times0{,}5=4$ points. Le total vaut $\res{13/20}$. Reprendre le cours de fractions, refaire les questions ratées sans regarder la solution, puis contrôler sur des exercices nouveaux cible la difficulté identifiée. Ce bilan guide la suite sans résumer toutes les compétences à une note.
\end{corrige}\end{blocLivre}
